\section{An exact finite expansion on one cell}
\label{sec:cell}
\label{sec:cell-expansion}

The identities in this section concern a finite cloud.  In particular, they
do not assume that large components, recollisions, or long encounters are
unlikely.  Their signed auxiliary trajectories are distinguished throughout
from the trajectory of the full microscopic cloud. The use of autonomous
physical trajectory components and signed geometric overlap constraints
has its antecedent in~\cite{BGSSCluster2022}. For smooth potentials,
the identities below retain entire interacting components and their
potential energy.

\subsection{Contact components and kernels}

Fix a cell $I=[a,b]$, put $\delta=b-a$, and let $S$ be a finite label set.
For $A\subset S$ write
\[
 X_i^A(u;Z_A),\quad V_i^A(u;Z_A),\qquad i\in A,\quad 0\leq u\leq\delta,
\]
for the isolated Hamiltonian flow of precisely the labels in $A$, started
from $Z_A$ at $a$.  The empty flow has its evident meaning.  Statements about
these flows are made outside their null sets of undefined initial data.
The flow and its measurability are supplied by Proposition~\ref{pre:flow}.

Let $\pi_0(S)=\pi_0(Z_S)$ be the partition into connected components of
the graph with edges $|x_i-x_j|\leq\varepsilon$.  A subset is
\emph{initially admissible} if it is a union of blocks of $\pi_0(S)$.
Let $\Pi_I(A)$ be the connected-component partition of the graph on $A$
with an edge $ij$ whenever
\[
 \min_{0\leq u\leq\delta}|X_i^A(u)-X_j^A(u)|\leq\varepsilon.
\]
Thus a contact includes a force-free passage or a tangency.  Set
\[
 c_I(A)=\mathbf 1_{\{\Pi_I(A)=\{A\}\}}\quad(A\ne\varnothing),
 \qquad
 c_I^R(A)=\mathbf 1_{\{\text{every block of }\Pi_I(A)\text{ meets }R\}},
 \quad R\subset A.
\]
The latter indicator is one for $A=R=\varnothing$.  It is zero for
nonempty $A$ and empty $R$.

If $A$ and $B$ are disjoint, their two trajectories in the following
definition are computed \emph{separately}:
\[
 \tau_I(A,B)=\mathbf 1_{\{\exists u\in[0,\delta],\ i\in A,\ j\in B:
             |X_i^A(u)-X_j^B(u)|\leq\varepsilon\}},
 \qquad \alpha_I(A,B)=1-\tau_I(A,B).
\]
For a collection $\mathcal B$ of disjoint force groups, use
\[
 \tau_I(\mathcal B,C)=1-\prod_{B\in\mathcal B}\alpha_I(B,C).
\]
The union of groups in $\mathcal B$ is not assigned a new interacting
flow in this notation.  Its paths remain the indicated separate paths.
For a prescribed collection of paths $\mathcal X$ and a remaining cloud
$V$, define
\begin{equation}
 \Xi_I(\mathcal X;V)
 =\prod_{C\in\Pi_I(V)}\bigl(1-\tau_I(\mathcal X,C)\bigr).
 \label{eq:cell-Xi}
\end{equation}
The components on the right follow their isolated flows, which coincide
with their restrictions of the $V$-flow.  Empty products are one.

For $R\subset A$ the endpoint isolation indicator is
\[
 e_I^R(A)=\prod_{i\in R}\prod_{j\in A\setminus R}
       \mathbf 1_{\{|X_i^A(\delta)-X_j^A(\delta)|>\varepsilon\}}.
\]
It imposes no separation between two labels of $R$.  Introduce the
positive measure kernel
\begin{equation}
 \mathsf H^I_{R,A}(Z_A;\dd Y_R)
   =e_I^R(A)\,
     \delta_{(X_i^A(\delta),V_i^A(\delta))_{i\in R}}(\dd Y_R).
 \label{eq:cell-H}
\end{equation}
The Dirac measure notation avoids introducing a density for a
lower-dimensional transport kernel.

The only flow fact needed for the algebra is the following equivalence.
If separate groups avoid one another throughout $I$, their separate
flows are the restrictions of the flow of their union.  Conversely,
different blocks of $\Pi_I(S)$ avoid one another, and their restrictions
are isolated flows.  Indeed cross forces vanish in the first assertion,
so uniqueness proves it; the second follows from the definition of the
contact graph and the same uniqueness argument.

\begin{definition}[A cell diagram]
\label{def:cell-diagram}
For $R\subset S$, with $R\ne\varnothing$, a cell diagram on $(R,S)$ is
\[
 \gamma=(A;\Gamma_1,\ldots,\Gamma_\ell).
\]
Here $R\subset A\subset S$; the nonempty sets in all the $\Gamma_g$
together partition $S\setminus A$; and each $\Gamma_g$ is a nonempty
unordered collection of these sets.  Put $\Gamma_0=\{A\}$.
Different members of the diagram are different force groups.  Let
$\iota_0(\gamma)$ be the indicator that every force group in the diagram
is a union of blocks of $\pi_0(S)$.  Equivalently, all pairs in different
force groups have distance greater than $\varepsilon$ at $a$.
Define
\begin{align}
 w_I(\gamma;Z_S)
  ={}&\iota_0(\gamma)c_I^R(A)
       \prod_{g=1}^{\ell}
       \left[
        \prod_{B\in\Gamma_g}c_I(B)\tau_I(\Gamma_{g-1},B)
        \prod_{\{B,C\}\subset\Gamma_g}\alpha_I(B,C)
       \right],                                      \label{eq:cell-weight}\\
 \sigma(\gamma)={}&(-1)^{\sum_{g=1}^{\ell}|\Gamma_g|}.
                                                        \label{eq:cell-sign}
\end{align}
Only the root group $A$ uses a joint flow containing roots.  In
particular, a contact between different generations is a condition on
separate trajectories, not a force exerted between their groups.
The empty list of generations is allowed when $S=A$.
\end{definition}

Define the signed kernel
\begin{equation}
 \mathsf K^I_{R,S}(Z_S;\dd Y_R)
  =\sum_{\gamma\text{ on }(R,S)}
       \sigma(\gamma)w_I(\gamma;Z_S)
       \mathsf H^I_{R,A}(Z_A;\dd Y_R).
 \label{eq:cell-K}
\end{equation}
The sum is finite.  Set $\mathsf K^I_{\varnothing,\varnothing}=1$ and
$\mathsf K^I_{\varnothing,S}=0$ for $S\ne\varnothing$.
For a function or density $g$ on the input coordinates define
$\mathcal T^I_{R,S}g$ as a signed measure by
\begin{equation}
 \langle\psi,\mathcal T^I_{R,S}g\rangle
 =\int g(Z_S)\int\psi(Y_R)
                   \mathsf K^I_{R,S}(Z_S;\dd Y_R)\,\dd Z_S.
 \label{eq:cell-T}
\end{equation}

\begin{remark}[One additional label]
\label{rem:cell-single}
If an extra label $e$ is initially separate from all roots, its two
diagram terms are
\[
 c_I^R(R\cup\{e\})\mathsf H^I_{R,R\cup\{e\}}
 \; -\;
 \tau_I(R,\{e\})\mathsf H^I_{R,R}.
\]
The first follows the full force group and the second follows the
root group and a free extra particle separately.  If $e$ belongs to
an initial root polymer, it is mandatory and only the full-group term
is present.  No extra family-selection multiplicity is introduced
by this identity.
\end{remark}

\subsection{Boolean inversion and the generation identity}

\begin{proposition}[Exact cell expansion]
\label{cell:expansion}
Let $R\subset S$ and let $A_0$ be the union of the blocks of
$\pi_0(S)$ meeting $R$.  For $R\ne\varnothing$, let $\mathcal P$ be
the collection of the other blocks of $\pi_0(S)$.  Then
\begin{equation}
 \mathsf H^I_{R,S}
 =\sum_{\mathcal U\subset\mathcal P}
      \mathsf K^I_{R,A_0\cup\bigcup\mathcal U},
 \label{eq:cell-subset-sum}
\end{equation}
and, as an equality of signed measures,
\begin{equation}
 \mathsf K^I_{R,A_0\cup\bigcup\mathcal U}
  =\sum_{\mathcal V\subset\mathcal U}
       (-1)^{|\mathcal U|-|\mathcal V|}
       \mathsf H^I_{R,A_0\cup\bigcup\mathcal V}.
 \label{eq:cell-boolean}
\end{equation}
Every coefficient is evaluated on its displayed input coordinates.
If $m=|S\setminus R|$, the number of diagrams in
\eqref{eq:cell-K} is at most $(C(1+m))^{Cm}$, with a universal $C$.
There is no partition of the roots in this count.
\end{proposition}

\begin{proof}
First decompose the full flow according to its unique union of contact
components meeting $R$.  The flow observation preceding
Definition~\ref{def:cell-diagram} gives
\begin{equation}
 \mathsf H^I_{R,S}
 =\sum_{\substack{R\subset A\subset S\\
                   A\text{ union of blocks of }\pi_0(S)}}
       c_I^R(A)\mathsf H^I_{R,A}\Xi_I(\{A\};S\setminus A).
 \label{eq:cell-root-decomposition}
\end{equation}
Precisely one term has its two component indicators nonzero: the actual
union of rooted contact components.  Endpoint isolation of $R$ against
the rest of $S$ can then be checked inside $A$, because $A$ and its
complement avoid throughout the cell.

Here is the finite identity that expands the last factor.  Fix
$\mathcal X$, and consider a collection $\mathcal B$ of disjoint
nonempty subsets of $V$, each a union of the fixed initial polymers.
Put $U_{\mathcal B}=\bigcup\mathcal B$ and
\[
 d_I(\mathcal B)
   =\prod_{B\in\mathcal B}c_I(B)
      \prod_{\{B,C\}\subset\mathcal B}\alpha_I(B,C).
\]
Then
\begin{equation}
 \Xi_I(\mathcal X;V)
 =\sum_{\mathcal B}
       (-1)^{|\mathcal B|}d_I(\mathcal B)
       \prod_{B\in\mathcal B}\tau_I(\mathcal X,B)
       \Xi_I(\mathcal B;V\setminus U_{\mathcal B}).
 \label{eq:cell-generation-recursion}
\end{equation}
The term with $\mathcal B=\varnothing$ equals one.  To prove
\eqref{eq:cell-generation-recursion}, expand the product
\eqref{eq:cell-Xi} over subsets of its actual components.  A given
collection $\mathcal B$ consists of components of $V$ exactly when
its members are connected, mutually avoid, and avoid the remaining
$V\setminus U_{\mathcal B}$ flow.  These conditions are respectively
$\prod c_I(B)$, the pair product in $d_I$, and the last $\Xi_I$ in
\eqref{eq:cell-generation-recursion}.  On their support the required
flows agree.  This proves the formula, including all signs.

Apply \eqref{eq:cell-generation-recursion} repeatedly to
\eqref{eq:cell-root-decomposition}.  Every nonempty selection strictly
decreases the number of remaining labels, so the process terminates.
At termination its coefficient is exactly
\eqref{eq:cell-weight}--\eqref{eq:cell-sign}; the unselected labels
have no remaining restriction.  Grouping by the union of selected
initial polymers proves \eqref{eq:cell-subset-sum}.  The initial
partition of such a union is the restriction of $\pi_0(S)$, so its
coefficient is exactly the kernel defined in \eqref{eq:cell-K}, not a
kernel depending on unselected coordinates.  Finite Boolean inversion
proves \eqref{eq:cell-boolean}.

There are at most $2^m$ choices of $A\setminus R$, at most $m^m$
partitions of the remaining labels, and at most $m^m$ assignments of
their blocks to successive nonempty generations.  This proves the
count, with the cases $m=0,1$ understood separately.
\end{proof}

\subsection{Expectations and normalization}

For any finite random cloud let its scaled factorial measure with
external isolation be denoted by $\widetilde F_R$.  Its weak definition
is the usual factorial sum with the additional indicator that every
label outside the displayed set has distance greater than
$\varepsilon$ from it.  This definition permits internal contacts.
For the actual evolved initial ensemble, all exponential moments of
the total particle number are finite.  Indeed its initial probability
of $N=n$ is bounded by $\mu^n/n!$, and $N$ is conserved.

\begin{corollary}[Factorial normalization]
\label{cor:cell-factorial}
For the actual microscopic cloud and $r=|R|$, the cell identity is
written after relabelling the ordered roots as $R=\{1,\ldots,r\}$:
\begin{equation}
 \widetilde F_R(b)
  =\sum_{m\geq0}\frac{\mu^m}{m!}
       \mathcal T^I_{R,R\sqcup\{r+1,\ldots,r+m\}}
       \widetilde F_{r+m}(a).
 \label{eq:cell-factorial}
\end{equation}
The right side is interpreted with each finite Boolean coefficient
grouped before summing over $m$.  Equivalently, one may first restrict
the ensemble to $N\leq M$, apply the finite identity, and then let
$M\to\infty$.  No assertion of absolute summability of the
\emph{ungrouped} generation series is made here.
\end{corollary}

\begin{proof}
Apply \eqref{eq:cell-subset-sum} to each ordered tuple of roots in a
realization of the cloud at $a$.  A selected union of its initial
polymers is exactly an externally isolated displayed subset at $a$.
Consequently the expectation of each selected union uses
$\widetilde F_{r+m}(a)$.  Its $m$ extra labels were unordered, whereas
the factorial sum orders them, giving the divisor $m!$.  The scaling
of the original $r$ roots is $\mu^{-r}$, and that of the new factorial
measure is $\mu^{-(r+m)}$, giving $\mu^m$.

For a bounded test $\psi$, Boolean inversion bounds a coefficient
in total variation by $2^{|\mathcal U|}\|\psi\|_\infty$.  Summing this
over all selected polymer subsets in a fixed cloud costs at most
$3^N\|\psi\|_\infty$, and the root tuple count costs at most $N^r$.
The integrable bound $\mu^{-r}N^r3^N\|\psi\|_\infty$ justifies the
stated grouped limit.  Expanding every large-order coefficient into
all its diagrams before taking absolute values would give a different,
unproved convergence assertion and is unnecessary.
\end{proof}

\subsection{Factorization on all displayed trajectories}

\begin{proposition}[Diagramwise factorization]
\label{cell:factorization}
Split the labels into $S=S_1\sqcup S_2$ and the roots into
$R=R_1\sqcup R_2$, where $R_i\subset S_i$.  Let $\gamma_i$ be a
diagram on $(R_i,S_i)$.  Define $\eta_I(\gamma_1,\gamma_2)$ to be the
indicator that every displayed particle trajectory of $\gamma_1$
has distance greater than $\varepsilon$ throughout $I$ from every
displayed trajectory of $\gamma_2$.  This includes detached groups
and includes the initial and terminal times.
For a single global diagram $\gamma$, define
$\eta_I(\gamma;S_1,S_2)$ by the same condition between its labelled
paths in $S_1$ and $S_2$, using their prescribed force-group flows.

On this indicator, merge the two root groups, and take the union of
the two batches at each generation, padding a finished list by empty
batches.  This gives a diagram on $(R,S)$.  Conversely, every nonzero
diagram term satisfying $\eta_I(\gamma;S_1,S_2)=1$ splits uniquely in
this manner.  Its detached groups necessarily belong to one $S_i$,
and its root flow splits over its two intersections with $S_i$.
Under this bijection the signed transport measures, with this
indicator imposed, agree with the tensor product of the two
diagram transport measures:
\begin{equation}
 \sigma(\gamma)w_I(\gamma)\mathsf H^I_{R,A}\,\eta_I
  =\eta_I\bigl[
       \sigma(\gamma_1)w_I(\gamma_1)\mathsf H^I_{R_1,A_1}
       \otimes
       \sigma(\gamma_2)w_I(\gamma_2)\mathsf H^I_{R_2,A_2}
             \bigr].
 \label{eq:cell-factorization}
\end{equation}
Summation over the bijection preserves this equality.  Factorial
weights split as
\begin{equation}
 \binom{m_1+m_2}{m_1}
       \frac{\mu^{m_1+m_2}}{(m_1+m_2)!}
   =\frac{\mu^{m_1}}{m_1!}\frac{\mu^{m_2}}{m_2!}.
 \label{eq:cell-factorial-split}
\end{equation}
\end{proposition}

\begin{proof}
Cross avoidance makes the joint root flow equal to the two separate
root flows.  Its rooted-component and endpoint-isolation indicators
therefore split.  Initial-polymer admissibility also splits, since
all cross initial distances are greater than $\varepsilon$.
Within a merged generation, all additional cross avoidance factors
are one.  A detached component in $S_i$ can touch the preceding
generation only through that generation's $S_i$ part.  Hence its
touch indicator is exactly its original touch indicator.  The
number of detached groups is additive, so the signs multiply.
The output Dirac measures split with the root flow.  These facts
give \eqref{eq:cell-factorization}.

For the inverse map, a connected detached group cannot contain labels
from both families: a connected graph across that partition would
have a cross contact, contrary to $\eta_I=1$.  The same avoidance
splits the root group's flow by uniqueness.  A nonempty generation
in one family must have
a predecessor in that same family unless it is the first generation:
otherwise its required touch indicator vanishes.  There can therefore
be no gap followed by a later nonempty generation in one family.
Removing its terminal empty batches produces its unique local
generation list.  This proves the asserted bijection.  Formula
\eqref{eq:cell-factorial-split} is the ordinary ordered-label count.
\end{proof}

The avoidance condition in Proposition~\ref{cell:factorization} is
\emph{diagramwise}.  It is not a single condition on the flow of $S$.
For example, a third particle can deflect a prospective collision
partner away from a root in the full flow while a smaller subflow
still changes that root.  Such configurations are included in the
subset coefficient and cannot be removed using full-flow separation.
For more than two factors the proposition applies successively; its
indicator must include all displayed paths from all factors.  This
is the precise factorization used by the centered identities below.

\subsection{A deterministic stopped identity}

We give a stopped version before any probabilistic tail estimate.
Its stopping parameter $k$ counts all selected labels, including
the roots.  The threshold is allowed to include labels already
selected in later cells: in that case subtract that fixed offset
from the remaining threshold in the construction below.

Order components by the lexicographically least initial position in
each component; denote this position by $\kappa(C)\in\mathbb R^3$.
Different actual contact components have different such positions,
since particles with the same initial position belong to the same
initial polymer.  This is a measurable, label-equivariant deterministic
order.  The following elementary weighted product
identity is useful.  For $I_1,\ldots,I_n\in\{0,1\}$, positive integer
weights $w_1,\ldots,w_n$, and an integer $d>0$, let
$w(J)=\sum_{j\in J}w_j$.  Then
\begin{align}
 \prod_{j=1}^n(1-I_j)
 ={}&\sum_{\substack{J\subset\{1,\ldots,n\}\\w(J)<d}}
             (-1)^{|J|}\prod_{j\in J}I_j \nonumber\\
   &+\sum_{\substack{\varnothing\ne J\subset\{1,\ldots,n\}\\
                w(J\setminus\{\max J\})<d\leq w(J)}}
       (-1)^{|J|}\prod_{j\in J}I_j
                \prod_{j>\max J}(1-I_j).
 \label{eq:cell-stopped-product}
\end{align}
To verify it, expand the factors in their prescribed order, stopping
a branch at its first selected factor whose accumulated weight is
at least $d$.  Branches that finish without stopping are the first
sum, and stopped branches are the second.  The factors remaining
on a stopped branch lie in $[0,1]$.

Apply \eqref{eq:cell-stopped-product} to
$\Xi_I(\mathcal X;V)$, with the actual components of $V$ as its
ordered indices, $I_j=\tau_I(\mathcal X,C_j)$ and $w_j=|C_j|$.
For a selected collection $\mathcal B$, denote its last component
in this order by $B_*$ and set $V'=V\setminus U_{\mathcal B}$.
After extracting the selected components, the extra factor on
a stopped branch is
\begin{equation}
 \vartheta_I(\mathcal X,\mathcal B;V')
   =\Xi_I(\mathcal B;V')
      \prod_{\substack{C\in\Pi_I(V')\\
                   \kappa(C)>_{\mathrm{lex}}\kappa(B_*)}}
                  \bigl(1-\tau_I(\mathcal X,C)\bigr).
 \label{eq:cell-stop-theta}
\end{equation}
It belongs to $[0,1]$.  The condition that the selected components
are actual components of $V$ is represented, exactly as in
\eqref{eq:cell-generation-recursion}, by $d_I(\mathcal B)$ and the
first factor in \eqref{eq:cell-stop-theta}.  On that condition the
order of the components of $V$ is the merged order of
$\mathcal B$ and $\Pi_I(V')$, which proves
\eqref{eq:cell-stop-theta}.  Completed branches with
$|U_{\mathcal B}|<d$ use the full extraction identity and continue
with $\Xi_I(\mathcal B;V')$.  A nonempty continued batch strictly
decreases the number of remaining labels.

\begin{proposition}[Stopped cell expansion and positive domination]
\label{cell:stopping}
For $k>|R|$, the finite-cloud expansion equals its complete diagrams
with fewer than $k$ selected labels plus a finite signed remainder.
Each remainder term is of one of the following types.
\begin{enumerate}
 \item Its root group $A$ already has $|A|\geq k$.  Its transport is
 $c_I^R(A)\mathsf H^I_{R,A}$ times the factor
 $\Xi_I(\{A\};S\setminus A)\in[0,1]$.
 \item Its completed generations and its final partial generation
 select a union $U$ with $|U|\geq k$.  Immediately before the last
 selected component $B_*$, the number of selected labels is less
 than $k$.  All displayed groups satisfy the connectivity,
 initial-polymer, same-generation avoidance, and preceding-generation
 touch conditions of Definition~\ref{def:cell-diagram}.  Its weight
 is the corresponding sign and indicators times a factor
 $\vartheta_I$ of the form \eqref{eq:cell-stop-theta}, lying in $[0,1]$.
\end{enumerate}
In the second case the overshoot is by at most one additional
component.  In particular, if every component of every trial
subsystem formed from initial polymers has size at most $D$, then
$|A|\leq |R|D$, and a stopped union has size at most
$\max\{|R|D,k-1+D\}$.

Replacing the sign by $+1$ and dropping the last factor gives a
pointwise positive domination of the total variation of the
remainder.  Its expected version is a sum of these positive
contact expressions integrated against the initial isolated
factorial measures of the selected unions, with the factors
$\mu^m/m!$ of \eqref{eq:cell-factorial}.  This is an identity and
a domination statement, not a claim that this sum tends to zero.
\end{proposition}

\begin{proof}
Use \eqref{eq:cell-root-decomposition}.  Stop a summand immediately
if its root group reaches the threshold.  Otherwise apply the
weighted product expansion with $d=k-|A|$, extract its selected
components, and continue its nonempty completed selections by the
same rule with the updated remaining threshold.  Empty selections
terminate a complete diagram.  The explicit finite identities
\eqref{eq:cell-stopped-product} and
\eqref{eq:cell-stop-theta} prove the resulting equality at every
stage.  A diagram of total size less than $k$ never triggers the
stopping rule and appears with exactly its original coefficient.
Every other branch stops in one of the two stated ways.

Weights are positive before their displayed inclusion-exclusion
sign, and the unexpanded factors are in $[0,1]$, proving the
variation bound.  All selected groups are unions of the original
initial polymers.  Their union is consequently externally isolated
in the original cloud at the initial cut; removing the last factor
does not remove this initial-polymer selection condition.  Taking
factorial expectations therefore gives the asserted input measures
and normalization.  On the stated trial-component bound each
component of the root group contains a root and has at most $D$
labels, proving $|A|\leq |R|D$; the other size bound follows from
the one-component overshoot.

These expected finite stopped identities are well defined for
the grand ensemble without a tail theorem.  Prior to the last
selected component there are fewer than $k$ selected labels, and
therefore at most $k$ selected groups.  The root-large branch has
one group.  At fixed $k$ the number of stopped descriptions in
an $N$-particle cloud is bounded by $C_k^{N+1}$: assign each label
to one of at most $k+2$ group or remainder slots, and then order
the bounded number of selected groups and generations.  The
exponential moments of $N$, with its root-tuple polynomial factor,
justify expectation and the finite-$N$ limit.  This argument
establishes finiteness only; it supplies no useful uniform bound
as $\varepsilon\to0$.
\end{proof}

\subsection{Whole-cell histories}

\begin{definition}[A prescribed whole-cell history]
\label{cell:histories}
A finite whole-cell history is a composition of the diagram kernels
above on successive cells, together with any tensor products and
singleton-slot deletions or independent insertions specified by
the exact centered identities at layer cuts.  Its discrete
prescription records the following data:
\begin{enumerate}
 \item every formal label and its lifetime, with a fresh label
 for each independent input and no reuse of an independent source;
 \item on each occupied cell, the required root labels $R$, the
 full set $S$, its root force group, and its detached groups and
 generations;
 \item all the indicators in \eqref{eq:cell-weight} and the
 endpoint isolation indicators, including the initial polymer
 restrictions restored at every deleted singleton slot;
 \item the signs, factorial weights, and any deterministic
 stopping marks and unexpanded factors.
\end{enumerate}
On an ordinary hierarchy cell, labels in $R$ persist through its upper endpoint;
labels in $S\setminus R$ end their cell lifetime there.  In the
forward direction they were present throughout the cell, not
created at a first-contact time.  The root group is isolated from
its departing members at the upper endpoint by $e_I^R(A)$.
Detached groups have separate force systems; their geometric
contacts with other generations at that endpoint are unrestricted.
At the lower endpoint all force groups are unions of initial
polymers.  At a layer cut a nonsingleton polymer remains intact;
independent singleton inputs are inserted only with the full
polymer-pattern restrictions of the centered identity.

The size $p$ is the number of distinct formal label lifetimes, not
the number of their repeated appearances at cuts.  All transport
maps and indicators are defined on the full coordinates of these
lifetimes.  A history with a stopped factor can still depend on
the unexpanded microscopic cloud through its last factor in
$[0,1]$; its positive majorant drops that factor.
\end{definition}

There is also a positive version needed before a stopped cell.  If a
nonnegative observable of displayed labels, carrying their external
isolation indicator at the top of a cell, is evaluated along the
actual cloud, extract its actual union of
whole-cell components containing those labels, as in
\eqref{eq:cell-root-decomposition}.  Retain its root endpoint
isolation requirement and drop external avoidance against the
remaining cloud for an upper bound.  Repeat backwards over earlier
cells.  This produces only positive root-group histories, each
connected by contacts to the preceding required set, with whole
initial polymers at each cut.  The construction is an exact
decomposition followed by inequalities between nonnegative
functions.  Polynomial caps or probabilities for such histories
are not presumed in this definition; they are proved in
\Cref{close:witnesses}.

\subsection{The terminal boundary of a positive witness}

The positive histories just described retain the usual endpoint
isolation at every cut.  To detect a large component in the actual
cloud, we will also use a final block whose distinguished label may
still be interacting at the observation endpoint.  The next variant
keeps that entire final force group in the integral.  This boundary
condition is needed for the energy estimate as well as for the
witness construction in Section~\ref{sec:closure}.

\begin{definition}[An unisolated terminal witness block]
\label{def:cell-terminal-witness}
A terminal witness is permitted a different upper boundary condition
from an ordinary hierarchy cell.  For $R\subset A$ put
\begin{equation}
 \mathsf J^I_{R,A}(Z_A;\dd Y_R)
   =\delta_{(X_i^A(\delta),V_i^A(\delta))_{i\in R}}(\dd Y_R),
 \label{eq:cell-J}
\end{equation}
with no endpoint isolation factor.  Its exact positive root-component
decomposition is
\begin{equation}
 \mathsf J^I_{R,S}
  =\sum_{\substack{R\subset A\subset S\\
                   A\text{ union of blocks of }\pi_0(S)}}
      c_I^R(A)\mathsf J^I_{R,A}\Xi_I(\{A\};S\setminus A).
 \label{eq:cell-unisolated-root-decomposition}
\end{equation}
Indeed the unique actual union of rooted contact components is still
the unique nonzero summand, and its root trajectories agree with the
full ones.  This proves the identity without asserting endpoint
isolation.  Integrating a nonnegative test and dropping its last
avoidance factor gives the corresponding positive bound.

For a terminal witness block, all selected force groups end their
formal lifetimes together at the fixed upper endpoint $b$.  The
distinguished labels $R$---usually a single label---serve only for
factorial normalization and spatial integration anchors.  They do
not persist as a force subsystem after that endpoint.  All terminal
coordinates of every selected group are retained in the integration
and energy bookkeeping before any projection in
\eqref{eq:cell-J}.  In particular one keeps the full terminal group
Hamiltonians; one does not split $R$ away from close companions.
Geometric isolation of $R$ from $A\setminus R$ at $b$ is neither
required nor inferred.

A selected block may also carry a measurable nonnegative trajectory
indicator $W_I(A;Z_A)\leq1$ specifying the witness.  Such an indicator
is evaluated on its isolated group flow.  If the witness was
identified at an earlier time inside the cell, keep that prefix
condition and continue the isolated group to the fixed endpoint
$b$ before using whole-cell coordinates.  This is a finite
prescription with an extra indicator, not a stopping-time change
of phase-volume variables.  Establishing that an actual bad event
admits a selected witness of this form belongs to
\Cref{close:witnesses}.

At every earlier cut the ordinary rules of
Definition~\ref{cell:histories} remain in force.  In particular, a
selected union is initially externally isolated and its backward
continuation across the preceding cell retains the usual endpoint
isolation indicator.  Only the final witness boundary uses
\eqref{eq:cell-J}.  These are formal operations in the witness
integral; the microscopic evolution is not stopped or modified.
\end{definition}

\subsection{Contact records and coordinate refinements}
\label{cell:contact-records}

The ownership rule at a centered cut needs a record of which contacts
connect new labels to the old array.  We specify the data here;
Lemma~\ref{hist:pinning} constructs the corresponding integration
charts and proves their Jacobian and counting bounds.

For each occupied cell and each original tensor factor, an
\emph{original contact record} is an ordered list of edges between
displayed labels.  An edge records its two labels and either their
relative displacement at an endpoint overlap or their contact time
and sphere normal.  Starting at the late endpoint, with all old
labels in one anchored coordinate group and the new labels in
separate groups, each edge joins two groups not previously joined.
The edge list connects every new label to the old array.
Endpoint overlaps are processed first; subsequent edges are ordered
by successive first backward contacts, with the prescribed
deterministic tie rule.  All edges stay within the original tensor
factor.  When there is exactly one added label across the cell and
it is separated from its own roots at the late endpoint, the record
retains its first backward contact with that factor and the old label
met there as its original parent.  These groupings record contact coordinates;
they do not replace any prescribed force group.

The algebraic generations and the geometric pinning forest
need not coincide.  A diagram can be partitioned measurably according
to the first contact of a floating group with \emph{any} already
displayed path, including a path from another tensor factor.  Such
a partition inserts indicators whose sum is one on the original
term, and leaves its force systems, signs, and original required
contact indicators unchanged.  A cross-factor geometric pin is
therefore permitted even if it is not an edge of its algebraic
generation.  All equalities above survive any such measurable
refinement.  The Jacobians and causal properties of the refinements
used for estimates are proved separately in
\Cref{hist:pinning}; they are not consequences of
Boolean inversion alone.
